\subsection{Hadamard row sectors with explicit resistant targets}\label{ian:hadamard-code-sectors}
Status: proved: a random balanced-target discrepancy bound for the row code of an arbitrary Hadamard matrix of even order (Lemma~\ref{ian:hadamard-random-balanced}); an explicit resistant target on tensor powers of the order-\(12\) Paley matrix, a code that is non-Sylvester in order only (Theorem~\ref{ian:hadamard-explicit-target}); efficient uniform one-query readout (Corollary~\ref{ian:hadamard-efficient-uniform}); and full-AA failure for the pure code laws and constant-weight mixtures, through property (P) (Corollary~\ref{ian:hadamard-full-aa}). The mechanism is not new. One-query identification of Hadamard rows is van Dam's (from literature: arXiv:quant-ph/0008059, Algorithmica \textbf{34} (2002), 413--428, Theorem~1 and Corollary~1); van Dam already treats Paley matrices (Section~5 of that paper), with an efficient implementation in van Dam--Hallgren (from literature: arXiv:quant-ph/0011067). By Proposition~\ref{ian:hadamard-character-encoding} the tensor code is a punctured character code with a non-uniform product label prior, and the target is the analogue of the block construction of Subsection~\ref{ian:hard-sector-instances} on a \(12\)-row base (\(\beta=5/6\) instead of \(1/2\)). New here: the random balanced-target bound and the hard-sector packaging of a code whose order is not a power of two (Remark~\ref{ian:hadamard-literature-comparison}).

Throughout, \(n=N\), logarithms with subscript \(2\) are binary, and \(\ln\) denotes the natural logarithm. Use the coordinate oracle
\[
 O_x|j,b\rangle
 =|j,b\oplus(1-x_j)/2\rangle .
\]
The classes \(\mathcal D_t,\mathcal F_t\), and hardness parameters follow
Definition~\ref{ian:hard-sector-definition} in
\path{ian/13_hard_sector_definitions_forward.tex}.
In particular, covariance tolerance \(\gamma\) for \(\mathcal F_t\) is equivalent to tolerance \(\gamma/2\) for \(\mathcal D_t\).

\begin{definition}[Hadamard row laws and acceptance functions]
\label{ian:hadamard-row-law}
Let \(H\in\{-1,1\}^{N\times N}\) satisfy \(HH^{\mathsf T}=NI\), and write \(c_a=(H_{aj})_j\). Let \(A\) be uniform on the row indices. The unsigned row law is \(\nu_H=\mathcal L(c_A)\). The signed row law is \(\nu_H^\pm=\mathcal L(\sigma c_A)\), where \(\sigma\) is an independent uniform sign.

For a public target \(g:\{1,\ldots,N\}\to[0,1]\), define on the entire cube
\[
 \widehat x_H(a)=\frac1N\sum_jH_{aj}x_j,
 \qquad
 Q_{H,g}(x)=\sum_a g(a)\widehat x_H(a)^2 .
 \tag{34.1}\label{ian:hadamard-global-acceptance}
\]
The target depends on \(A\), not on \(\sigma\).
\end{definition}

\begin{proposition}[One-query readout for an arbitrary Hadamard matrix]
\label{ian:hadamard-one-query}
The function \(Q_{H,g}\) belongs to \(\mathcal Q_{N,1}\), and
\[
 0\le Q_{H,g}(x)\le1
 \quad(x\in\{-1,1\}^N),
 \qquad
 Q_{H,g}(\sigma c_a)=g(a).
\]
The circuit identifies \(a\) exactly on either row law.
\end{proposition}
\begin{proof}
Prepare \(N^{-1/2}\sum_j|j\rangle\), use the answer register \(|-\rangle\) to make one phase query, and apply \(U_H=H/\sqrt N\). The resulting state is
\[
 \sum_a\left(\frac1N\sum_jH_{aj}x_j\right)|a\rangle .
\]
Measure \(a\), and accept with probability \(g(a)\). This gives
\eqref{ian:hadamard-global-acceptance}. Orthogonality implies
\[
 \sum_a\widehat x_H(a)^2
 =\frac1N\sum_jx_j^2=1.
\]
For \(x=\sigma c_a\), the state after \(U_H\) is \(\sigma|a\rangle\). The global sign does not affect the measurement.
\end{proof}

\begin{remark}[Query cost, gate cost, and uniformity]
\label{ian:hadamard-general-gates}
The proposition charges one coordinate query. It separately requires preparation of the uniform state, implementation of \(H/\sqrt N\), and evaluation of \(g\). A general public matrix can be compiled by finite-dimensional unitary decomposition, with polynomial cost in \(N\); this gives no polynomial-in-\(\log N\) implementation. Efficient evaluation of matrix entries alone does not supply an efficient basis change. Likewise, the random-target argument below supplies no efficient target evaluator. The explicit tensor family below has efficient preparation, basis change, and target evaluation.
\end{remark}

\begin{lemma}[Random balanced-target discrepancy]
\label{ian:hadamard-random-balanced}
Fix a Hadamard matrix \(H\) of even order \(N\). Choose \(g\) uniformly among Boolean targets taking the value \(1\) on exactly \(N/2\) rows. Put
\[
 M_{N,t}=N^{2^t-1}2^{2^t}.
\]
For every integer \(t\ge0\) and every \(\eta>0\),
\[
 \Pr_g\left\{
 \sup_{F\in\mathcal F_t}
 |\Cov_{\nu_H}(Q_{H,g},F)|>\eta
 \right\}
 \le 2(N+1)M_{N,t}e^{-2N\eta^2}.
 \tag{34.2}\label{ian:hadamard-random-tail}
\]
Consequently, with probability at least \(1-\delta\),
\[
 \sup_{F\in\mathcal F_t}|\Cov_{\nu_H}(Q_{H,g},F)|
 \le
 \sqrt{\frac{\ln M_{N,t}+
                 \ln(2(N+1)/\delta)}{2N}} .
\]
On the same event, the supremum over \(\mathcal D_t\) is at most half this bound.

In particular, for \(N\ge2^{32}\), with probability at least
\(1-e^{-N/512}\), simultaneously for every
\(t\le\lfloor\frac12\log_2N\rfloor\),
\[
 |\Cov_{\nu_H}(Q_{H,g},F)|\le\frac1{16}
 \quad(F\in\mathcal F_t),
 \qquad
 |\Cov_{\nu_H}(Q_{H,g},D)|\le\frac1{32}
 \quad(D\in\mathcal D_t).
\]
Thus
\[
 \operatorname{Hard}\left(
 Q_{H,g},\nu_H;t,\frac14,\frac1{16}\right).
\]
All assertions also hold for \(\nu_H^\pm\).
\end{lemma}
\begin{proof}
First let \(\varepsilon_a\) be independent uniform signs. For a fixed tree with Boolean signed leaf outputs, put \(f_a=F(c_a)\in\{-1,1\}\). Hoeffding's inequality gives
\[
 \Pr\left\{
 \left|\sum_a\varepsilon_af_a\right|>2N\eta
 \right\}
 \le2e^{-2N\eta^2}.
\]
Conditioning on \(\sum_a\varepsilon_a=0\) makes
\(g(a)=(1+\varepsilon_a)/2\) a uniformly random balanced target. The conditioning event has probability
\[
 2^{-N}\binom N{N/2}\ge\frac1{N+1},
\]
because the central binomial coefficient is maximal among the \(N+1\) coefficients. Under this conditioning,
\[
 \Cov_{\nu_H}(Q_{H,g},F)
 =\frac1{2N}\sum_a\varepsilon_af_a.
\]
The conditional failure probability is therefore at most
\(2(N+1)e^{-2N\eta^2}\).

Lemma~\ref{ian:hard-sector-tree-closure} counts at most \(M_{N,t}\) Boolean signed trees. It also reduces arbitrary signed leaf outputs to these extremal assignments. A union bound proves \eqref{ian:hadamard-random-tail}, including adaptive trees whose strategy is chosen after the public target is known. The relation \(F=2D-1\) gives the factor two for \(\mathcal D_t\).

For the numerical specialization, take \(\eta=1/16\) and the largest allowed depth. Then
\[
 \ln M_{N,t}\le\sqrt N\ln(2N).
\]
For \(N\ge2^{32}\),
\[
 \sqrt N\ln(2N)\le N/256,
 \qquad
 \ln(2(N+1))\le N/512.
\]
The logarithm of the failure bound is at most
\[
 N/256+N/512-N/128=-N/512.
\]
Shallower classes are contained in the largest class, and balancedness gives variance \(1/4\).

Finally, for each fixed \(\sigma\), the function \(x\mapsto F(\sigma x)\) is a depth-\(t\) coordinate tree. Since the target mean is independent of \(\sigma\), averaging the two covariance bounds proves the signed assertion.
\end{proof}

\begin{definition}[An explicit Paley tensor family]
\label{ian:hadamard-paley-tensor}
Let \(R=\{\infty\}\cup\mathbb F_{11}\). Define the quadratic character by
\[
 \chi(0)=0,\qquad
 \chi(u)=1\quad(u\in\{1,3,4,5,9\}),\qquad
 \chi(u)=-1\quad(u\in\{2,6,7,8,10\}).
\]
Define the \(12\times12\) matrix \(B\), with rows and columns indexed by \(R\), by
\[
 B_{\infty b}=B_{a\infty}=1,\qquad
 B_{aa}=-1\quad(a\in\mathbb F_{11}),\qquad
 B_{ab}=\chi(b-a)\quad(a,b\in\mathbb F_{11},\ a\ne b).
\]
This matrix is equivalent to Paley construction I at \(q=11\) (from literature:
R.~E.~A.~C.~Paley, \emph{On Orthogonal Matrices},
Journal of Mathematics and Physics \textbf{12} (1933), 311--320,
DOI \texttt{10.1002/sapm1933121311}).
The needed properties are verified below.

For \(k\ge1\), put
\[
 N=12^k,\qquad H^{(k)}=B^{\otimes k},
 \qquad
 H^{(k)}_{a,j}=\prod_{\ell=1}^kB_{a_\ell j_\ell}
 \quad(a,j\in R^k).
\]
Let \(A\) be uniform on \(R^k\), and let \(\nu_k\) be the uniform law on the rows of \(H^{(k)}\). Define
\[
 f_0(\infty)=-1,\qquad f_0(a)=1\quad(a\in\mathbb F_{11}),
 \qquad
 Z_k(a)=\prod_{\ell=1}^kf_0(a_\ell),
 \qquad
 g_k(a)=\frac{1-Z_k(a)}2 .
\]
Thus \(g_k(a)\) is the parity of the number of coordinates equal to \(\infty\).

The explicit global acceptance function is
\[
 Q_k(x)=
 \sum_{\substack{a\in R^k\\
       \sum_\ell\mathbf1_{\{a_\ell=\infty\}}\equiv1\pmod2}}
 \left(
 \frac1{12^k}\sum_{j\in R^k}
 x_j\prod_{\ell=1}^kB_{a_\ell j_\ell}
 \right)^2 .
 \tag{34.3}
\]
\end{definition}

\begin{lemma}[Paley orthogonality and a local correlation gap]
\label{ian:hadamard-paley-local-gap}
One has \(BB^{\mathsf T}=12I\). Every column of \(B\) has an even number of negative entries. For every finite list \(j_1,\ldots,j_s\in R\), including the empty list,
\[
 \left|
 \frac1{12}\sum_{a\in R}
 f_0(a)\prod_{r=1}^sB_{a j_r}
 \right|
 \le\frac56.
 \tag{34.4}\label{ian:hadamard-local-gap}
\]
Moreover, \(H^{(k)}\) is not equivalent to a Sylvester matrix. This is a statement about the order only: by Proposition~\ref{ian:hadamard-character-encoding}, the row code is a punctured character code.
\end{lemma}
\begin{proof}
For distinct \(a,b\in\mathbb F_{11}\),
\[
 \sum_z\chi(z-a)\chi(z-b)=-1.
\]
Here is an elementary verification. For nonzero \(d\), the equation
\(v^2=u^2-d\) has \(10\) solutions: choose \(u-v\ne0\), and then \(u+v=d/(u-v)\). Counting instead by \(u\) gives
\[
 11+\sum_u\chi(u^2-d)=10.
\]
Completing the square in \((z-a)(z-b)\) proves the asserted identity.

Let \(C_{ab}=\chi(b-a)\) on the finite indices, and let \(J\) be the all-one matrix of order \(11\). Since \(\chi(-1)=-1\), the preceding identity gives
\[
 C^{\mathsf T}=-C,\qquad C\mathbf1=0,\qquad
 CC^{\mathsf T}=11I-J.
\]
The finite block of \(B\) is \(C-I\). Its row sums are \(-1\), and
\[
 J+(C-I)(C^{\mathsf T}-I)=12I.
\]
Together with the all-one border, this proves \(BB^{\mathsf T}=12I\).

The \(\infty\)-column has no negative entries. Every other column is orthogonal to the all-one \(\infty\)-column, and therefore has six negative entries. Consequently the function
\(v(a)=\prod_rB_{a j_r}\) satisfies \(\prod_{a\in R}v(a)=1\).
Since \(\prod_{a\in R}f_0(a)=-1\), the function \(f_0v\) has an odd number of negative entries. That number lies between \(1\) and \(11\), so its average has absolute value at most \(10/12\). This proves \eqref{ian:hadamard-local-gap}.

Tensor products preserve Hadamard orthogonality. Finally, \(12^k\) is not a power of two, whereas a Sylvester matrix has power-of-two order. Row and column permutations and sign changes preserve the order. More generally, a square character table with all entries in \(\{-1,1\}\) must come from an elementary abelian \(2\)-group, whose order is a power of two.
\end{proof}

\begin{theorem}[An explicit resistant target on a code that is non-Sylvester in order only]
\label{ian:hadamard-explicit-target}
Put \(\beta=5/6\). For every \(k\ge1\),
\[
 \E_{\nu_k}Q_k=\frac{1-\beta^k}{2},
 \qquad
 \Var_{\nu_k}Q_k=\frac{1-\beta^{2k}}4,
\]
and every adaptive coordinate tree satisfies
\[
 |\Cov_{\nu_k}(Q_k,F)|\le2^t\beta^k
 \quad(F\in\mathcal F_t).
 \tag{34.5}\label{ian:hadamard-explicit-covariance}
\]
For \(k\ge64\), put
\[
 t_k=\left\lfloor\frac{\log_2N}{32}\right\rfloor,
 \qquad N=12^k.
\]
Then
\[
 \Var_{\nu_k}Q_k\ge\frac3{16},
 \qquad
 \sup_{F\in\mathcal F_{t_k}}
 |\Cov_{\nu_k}(Q_k,F)|\le\frac1{256}.
\]
In particular,
\[
 Q_k\in\mathcal Q_{N,1},
 \qquad
 \operatorname{Hard}\left(
 Q_k,\nu_k;t_k,\frac18,\frac1{32}\right).
 \tag{34.6}\label{ian:hadamard-explicit-hard-sector}
\]
The same conclusions hold for the uniform signed row law.
\end{theorem}
\begin{proof}
On the code, Proposition~\ref{ian:hadamard-one-query} gives
\(Q_k(c_A)=g_k(A)\). Independence of the \(A_\ell\), together with
\(\E f_0(A_\ell)=\beta\), gives the mean and variance formulas.

For any set \(S\) of physical coordinates, independence and
\eqref{ian:hadamard-local-gap} give
\[
 \left|
 \E_A\left[Z_k(A)\prod_{j\in S}H^{(k)}_{A,j}\right]
 \right|
 =
 \prod_{\ell=1}^k
 \left|
 \E_{A_\ell}\left[
 f_0(A_\ell)\prod_{j\in S}B_{A_\ell j_\ell}
 \right]\right|
 \le\beta^k.
\]
Writing \(p_k=(1-\beta^k)/2\), it follows that
\[
 \left|
 \E_{\nu_k}\left[(Q_k-p_k)\prod_{j\in S}X_j\right]
 \right|
 \le\frac12\left(\beta^k+
       \beta^k\left|\E_{\nu_k}\prod_{j\in S}X_j\right|\right)
 \le\beta^k.
 \tag{34.7}\label{ian:hadamard-centered-moments}
\]

Fix an adaptive tree. Along a particular leaf, its queried coordinates
\(j_1,\ldots,j_d\) and answers \(b_1,\ldots,b_d\) are fixed, with \(d\le t\).
The leaf event \(L\) has indicator
\[
 \mathbf1_L(X)=2^{-d}\prod_{r=1}^d(1+b_rX_{j_r}).
\]
Expanding this product and applying
\eqref{ian:hadamard-centered-moments} gives
\[
 \left|\E_{\nu_k}[(Q_k-p_k)\mathbf1_L]\right|\le\beta^k.
\]
Repeated coordinates cause no difficulty, since \(X_j^2=1\).
There are at most \(2^t\) leaves. Multiplication by their outputs in
\([-1,1]\), followed by summation, proves
\eqref{ian:hadamard-explicit-covariance}. This argument controls adaptive
strategies directly, without converting a fixed-set conditional bound.

Since \(\beta^4<1/2\), for \(k\ge4\) one has
\(\beta^{2k}<1/4\), giving variance at least \(3/16\).
Also \(N<2^{4k}\), so \(t_k\le k/8\), and
\[
 2^{t_k}\beta^k
 \le2^{k/8}2^{-k/4}
 =2^{-k/8}\le2^{-8}
 \quad(k\ge64).
\]
These estimates imply \eqref{ian:hadamard-explicit-hard-sector}.
The signed assertion follows by conditioning on the independent global
sign and pulling it through the tree, as in
Lemma~\ref{ian:hadamard-random-balanced}.

The coefficient \(1/32\) in \(t_k\) is conservative. Since
\(2^t\beta^k=2^{t-k\log_2(6/5)}\), the bound
\eqref{ian:hadamard-explicit-covariance} tends to zero along depths
\(t\le c\log_2N\) for every fixed
\(c<\log_2(6/5)/\log_2 12\approx0.073\).
\end{proof}

\begin{corollary}[Uniform implementation and finite precision]
\label{ian:hadamard-efficient-uniform}
The family \(H^{(k)},g_k,Q_k\) has a uniform description using the parameter \(k\) and fixed finite data. Its exact one-query circuit uses \(O(k)\) fixed gates acting on at most four qubits, apart from polynomial-in-\(k\) reversible address conversion. The target is evaluated in \(O(k)\) operations.

Over a fixed efficiently universal gate library, uniform approximations use
\(\operatorname{poly}(k,\log\delta^{-1})\) gates and one coordinate query, with a global acceptance function \(\widetilde Q_k\) satisfying
\[
 \sup_x|\widetilde Q_k(x)-Q_k(x)|\le\delta.
\]
For \(\delta\le1/1024\) and \(k\ge64\),
\[
 \operatorname{Hard}\left(
 \widetilde Q_k,\nu_k;t_k,\frac18,\frac1{32}\right).
\]
\end{corollary}
\begin{proof}
Let \(U=B/\sqrt{12}\), embedded as \(V=U\oplus I_4\) on four qubits.
Because the \(\infty\)-row of \(B\) is all one,
\(U^{\mathsf T}|\infty\rangle\) is the uniform state on \(R\).
Prepare \(k\) such states, make the single phase query, apply \(U^{\otimes k}\), and measure. The final classical computation counts occurrences of \(\infty\) modulo two.

Thus preparation and decoding use \(2k\) fixed constant-dimensional gates.
Their entries are computable algebraic numbers; a fixed decomposition into bounded-arity gates can be chosen once. Converting the \(k\) base-\(12\) digits to the usual binary coordinate address, and uncomputing the address workspace, costs \(\operatorname{poly}(k)\) gates and no oracle calls. The matrix entries and target are also computed directly from their digit formulas.

Approximating each constant block to operator-norm error
\(O(\delta/k)\) gives, by telescoping the circuit, uniform acceptance error at most \(\delta\). Assign any fixed acceptance value to unused encodings.

For completeness, if \(P,Q\in[0,1]\) and \(\E|P-Q|\le\delta\), then
\[
 |\Var P-\Var Q|\le2\delta,
 \qquad
 |\Cov(P,F)-\Cov(Q,F)|\le2\delta
 \quad(|F|\le1).
\]
The first inequality follows by minimizing \(\E(P-c)^2\) over
\(c\in[0,1]\) and using
\(|(P-c)^2-(Q-c)^2|\le2|P-Q|\).
For the second, center the test and use \(|F-\E F|\le2\).
Consequently,
\[
 \Var_{\nu_k}\widetilde Q_k
 \ge\frac3{16}-\frac2{1024}>\frac18,
 \qquad
 |\Cov_{\nu_k}(\widetilde Q_k,F)|
 \le\frac1{256}+\frac2{1024}<\frac1{32}.
\]
\end{proof}

\begin{proposition}[A constant-query character encoding]
\label{ian:hadamard-character-encoding}
Although \(H^{(k)}\) is not a Sylvester matrix, its unsigned promise oracle is equivalent, within a factor of \(11\) in coordinate queries, to a character oracle with a specified nonuniform label prior. The factor \(11\) holds for coherent queries: the eleven bit queries act on a common answer qubit, equivalently as phase queries with a common control. An implementation that uses fresh answer registers and then uncomputes them costs \(22\).
\end{proposition}
\begin{proof}
For \(a\in R^k\), define \(v(a)\in\mathbb F_2^{11k}\) by
\[
 v(a)_{\ell r}=\frac{1-B_{a_\ell r}}2
 \quad(r\in\mathbb F_{11}).
\]
The expanded character oracle is the Boolean string
\[
 Y_s=(-1)^{s\cdot v(a)}
 \quad(s\in\mathbb F_2^{11k}).
\]
Its length \(2^{11k}=N^{11/\log_2 12}\) is polynomial in \(N\).

For an original coordinate \(j\in R^k\), let \(s(j)\) have one selected bit \(r=j_\ell\) in block \(\ell\) when \(j_\ell\ne\infty\), and no selected bit otherwise. Then
\[
 H^{(k)}_{a,j}=Y_{s(j)}.
\]
One expanded-oracle query therefore simulates an original query; coherently, compute \(s(j)\) reversibly from \(j\), make the query, and uncompute \(s(j)\).

Conversely, for each \(r\in\mathbb F_{11}\), define \(j^{(r)}\in R^k\) by taking \(j^{(r)}_\ell=r\) when \(s_{\ell r}=1\), and \(\infty\) otherwise. Then
\[
 Y_s=\prod_{r\in\mathbb F_{11}}H^{(k)}_{a,j^{(r)}}.
\]
Eleven original coordinate queries compute this product. Coherently, compute the indices \(j^{(r)}\) reversibly from \(s\) and apply \(O_x\) at each of them to one common answer qubit. The answers accumulate as
\[
 b\longmapsto b\oplus\bigoplus_{r\in\mathbb F_{11}}\frac{1-H^{(k)}_{a,j^{(r)}}}2
 =b\oplus\frac{1-Y_s}2 .
\]
In the conjugate basis of the answer qubit these are eleven controlled phase queries with a common control, whose product is the controlled phase of \(Y_s\); no answer workspace remains, so eleven queries suffice. If each answer is instead written to a fresh register, the eleven answers must be uncomputed after the parity is formed, and the cost is \(22\) queries.

The map \(a\mapsto v(a)\) is injective because distinct Hadamard rows differ. Uniform \(a\) therefore induces the uniform prior on \(12^k\) specified vectors, rather than the uniform law on \(\mathbb F_2^{11k}\). This prior is the \(k\)-fold product of the uniform law on the twelve vectors \((v(a)_{r})_{r\in\mathbb F_{11}}\), \(a\in R\).
\end{proof}

\begin{remark}[Literature comparison and scope of the contribution]
\label{ian:hadamard-literature-comparison}
One-query identification of the rows of a Hadamard matrix is known (from literature: Wim van Dam, \emph{Quantum Algorithms for Weighing Matrices and Quadratic Residues}, Algorithmica \textbf{34}(4) (2002), 413--428, DOI \texttt{10.1007/s00453-002-0975-4}, arXiv:quant-ph/0008059, Theorem~1 and Corollary~1). Proposition~\ref{ian:hadamard-one-query} is its Hadamard case followed by a target readout. Van Dam already treats Paley matrices (Section~5 of the same paper, through shifted Legendre sequences), and an efficient quantum implementation for that case is given by van Dam and Hallgren (from literature: Wim van Dam and Sean Hallgren, \emph{Efficient Quantum Algorithms for Shifted Quadratic Character Problems}, arXiv:quant-ph/0011067). The Sylvester case is the Bernstein--Vazirani procedure (from literature: Ethan Bernstein and Umesh Vazirani, \emph{Quantum Complexity Theory}, SIAM Journal on Computing \textbf{26}(5), 1411--1473 (1997), DOI \texttt{10.1137/S0097539796300921}).

The tensor construction does not supply a new mechanism. By Proposition~\ref{ian:hadamard-character-encoding}, the code \(H^{(k)}\) is a punctured character code: its rows are characters of \(\mathbb F_2^{11k}\) restricted to the \(12^k\) points \(s(j)\), and the uniform row law becomes a non-uniform product prior on the character labels; simulating the full, unpunctured character oracle costs a factor of \(11\) in queries. The target \(g_k\) is the analogue, with base target \(f_0\), of the block construction of Subsection~\ref{ian:hard-sector-instances} (the target \(h\) of Lemma~\ref{ian:hard-sector-orbit-covariance} and Proposition~\ref{ian:hard-sector-character-instance} in \path{ian/17_hard_sector_instances.tex}): a parity over blocks of a fixed base target, here on a \(12\)-row base with base correlation \(\beta=5/6\) in \eqref{ian:hadamard-local-gap} instead of \(1/2\). The non-Sylvester property of Lemma~\ref{ian:hadamard-paley-local-gap} concerns the order only.

Identification hardness alone does not imply hardness of a selected decision bit: a target could equal one queried column. New here are the random balanced-target bound of Lemma~\ref{ian:hadamard-random-balanced}, valid for every Hadamard matrix of even order, and the hard-sector packaging of a code whose order is not a power of two (Theorem~\ref{ian:hadamard-explicit-target} and Corollaries~\ref{ian:hadamard-efficient-uniform} and~\ref{ian:hadamard-full-aa}). No separation from character encodings is claimed.
\end{remark}

\begin{corollary}[Full-AA failure for specified law families]
\label{ian:hadamard-full-aa}
Define a law family by \(\mu_{12^k}=\nu_k\), with arbitrary laws at other dimensions. This family fails full AA. More generally, the same conclusion holds whenever
\[
 \mu_{12^k}\ge w_0\nu_k
\]
for some fixed \(w_0>0\) along an unbounded sequence of \(k\).

For any supplied domination \(\mu_N\ge w_N\nu_k\),
\[
 \ell_{\mu_N}(t_k,Q_k)\ge\frac{3w_N}{32}.
 \tag{34.8}\label{ian:hadamard-mixture-loss}
\]
\end{corollary}
\begin{proof}
Equation~\eqref{ian:hadamard-explicit-hard-sector} has
\[
 T=1,\qquad a=\frac18,\qquad \gamma=\frac1{32}=\frac a4.
\]
Proposition~\ref{ian:hard-sector-forward} proves
\eqref{ian:hadamard-mixture-loss}.

The precise criterion needed is property (P) in
Remark~\ref{ian:hard-sector-polynomial-quantifiers}, in
\path{ian/15_hard_sector_sets_asymptotics.tex}:
there must exist fixed \(K,p,q\ge1\), chosen before the budget constants,
such that for every \(A,B\ge1\) there is a quantum-realizable hard sector with
\[
 r>8,\qquad
 t>\left\lceil A(T+r)^B\right\rceil,\qquad
 w^{-1}\le Kr^p,\qquad a^{-1}\le Kr^q,
 \qquad
 \operatorname{Hard}(Q,\nu;t,a,a/4).
\]
For the pure laws take \(K=p=q=1\), \(r=16\), and \(w=1\).
Since \(t_k\to\infty\), choose \(k\) with
\(t_k>\lceil A17^B\rceil\). All conditions hold.

For domination by fixed \(w_0\), take
\(r=\max\{16,w_0^{-1}\}\), again fixed before choosing \(k\).
The same argument proves (P) for smooth sectors. The converse direction of
Remark~\ref{ian:hard-sector-polynomial-quantifiers}, which applies
Lemma~\ref{ian:hard-sector-aa-exclusion} to these witnesses, then gives
full-AA failure.
\end{proof}

\begin{remark}[Mixture weights and amplification]
\label{ian:hadamard-mixtures-amplification}
For
\[
 \mu_N=(1-w_N)\rho_N+w_N\nu_k
\]
with arbitrary \(\rho_N\), the code law is a smooth sector of weight \(w_N\).
For the uniform-cube background, it is also the exact conditional law on the unsigned code \(C_k\), whose actual set mass is
\[
 \mu_N(C_k)=w_N+(1-w_N)N2^{-N}.
\]
This actual mass may replace \(w_N\) in the sector estimates.

For varying positive weights, a sufficient condition for these witnesses to establish (P) is
\[
 \text{for every }B\ge1,\qquad
 \sup_k
 \frac{t_k}{(1+\max\{16,w_N^{-1}\})^B}
 =\infty.
\]
Indeed, take \(r=\max\{16,w_N^{-1}\}\), \(K=p=q=1\), and choose the index after \(A,B\).
Constant weights satisfy this condition, as do weights with
\(w_N^{-1}=t_k^{o(1)}\).
A weight described only as inverse-polynomial in \(N\), or even as
\((\log N)^{-d}\) for fixed \(d>0\), does not make the proved logarithmic depth exceed every polynomial in the relevant inverse accuracy.

A single finite matrix supplies a finite-resource witness. The tensor family supplies unbounded hiding depth while retaining one quantum query, constant dispersion, and weight one under its own promise laws. This tensorization changes the law family; it does not increase the weight of a supplied mixture component. Repeating a decoder for error reduction changes \(T\) and supplies no replacement for the quantifiers in (P).
\end{remark}
